\section{Retained-edge planar tools}\label{sec:retained-edge-tools}

We prove that a row whose evaluation is nearly constant along many
lines agrees, on many incidences, with a constant row plus a multiple
of \(Jz\), where \(J(t,Z)=(-Z,t)\).  This is
\Cref{a06:row-fit}, used in both the finite-narrowness reduction and
the critical-rate argument.

The main step is planar rigidity for graphs with small restricted
sumsets: many original edges have their endpoints in two parallel
strips, of width equal to the mesh up to a subpower factor and with
subpower-bounded common slope.  Keeping original edges allows this
geometric conclusion to be lifted back to incidence mass.  We first
prove the planar statement, then apply it to the line projections
that arise from the row.

Throughout this section \(h=N^{-a}\), where \(a>0\) is fixed, and
\(K=h^{-o(1)}\) may increase by a fixed power from line to line.
For a bounded planar set \(U\), \(N_h(U)\) is its covering number by
squares of side \(h\).  Fixed-grid and separated-set versions differ
only by dimensional constants.  Constants indexed by a fixed
expression length may depend on that length, but remain subpower in
\(h^{-1}\).  The final passage to strip widths \(h^{1-o(1)}\) uses
a slow diagonal and may introduce a new subpower factor.

\subsection{Graph refinement and common endpoint sets}

We begin with the graph form of Balog--Szemer\'edi--Gowers that records
the edges surviving the sumset reduction. The three-path argument
follows \citet[Lemma~4.2 and the proof of Theorem~4.1]{SudakovSzemerediVu2005}; compare also
\citet[Theorem~2.29 and Exercise~6.4.10]{TaoVu2006} for the graph
form and retention of original edges. We give the proof with the
polynomial dependence needed after mesh rounding.

\begin{lemma}[Three paths retain graph edges]\label{re:graph-bsg}
Let \(A,B\) be nonempty finite subsets of an abelian group, with
\(a_0=|A|\), \(b_0=|B|\), and let \(G\subset A\times B\) have at least
\(\theta a_0b_0\) edges, \(0<\theta\le1\).  Put
\(C=\{x+y:(x,y)\in G\}\).  There are \(A'\subset A\), \(B'\subset B\)
such that
\begin{align*}
 |A'|&\ge c\theta a_0,\qquad |B'|\ge c\theta b_0,\\
 |G\cap(A'\times B')|&\ge c\theta^2a_0b_0,\qquad
 |A'+B'|\le C_0\theta^{-5}\frac{|C|^3}{a_0b_0},
\end{align*}
where \(c,C_0>0\) are absolute constants.
\end{lemma}
\begin{proof}
We first choose endpoint sets with many three-edge walks between every
pair.  The sums along those walks will bound the full sumset, while
degree counting will retain many edges of the original graph.
Remove vertices of \(A\) of degree below \(\theta b_0/2\), retaining
a set \(A_0\) and at least \(\theta a_0b_0/2\) edges.  An ordered pair
in \(A_0^2\) is bad if its common neighborhood has fewer than
\(\tau b_0\) vertices, where \(\tau=\theta^3/512\).  For uniform
\(v_0\in B\) let \(U=N(v_0)\cap A_0\).  Degree counting and Jensen give
\[
 \E|U|\ge\theta a_0/2,\qquad
 \E|U|^2\ge\theta^2a_0^2/4,\qquad
 \E\operatorname{Bad}(U)\le\tau a_0^2,
\]
the last bound because a bad pair has probability below \(\tau\) of
lying in \(U^2\).  For one \(U\), therefore,
\[
 |U|^2-(32/\theta)\operatorname{Bad}(U)
 \ge 3\theta^2a_0^2/16.
\]
It follows that \(|U|\ge\theta a_0/(2\sqrt2)\) and
\(\operatorname{Bad}(U)\le\theta|U|^2/32\).  Delete vertices having
more than \(\theta|U|/16\) bad partners.  At most half of \(U\) is
deleted; call the remainder \(A'\).  Define
\[
 B'=\{v\in B:|N(v)\cap A'|\ge\theta|A'|/4\}.
\]
Every member of \(A'\) retains degree at least \(\theta b_0/2\) in
the original graph.  Edges to \(B\setminus B'\) number at most
\(\theta|A'|b_0/4\), so
\[
 |G\cap(A'\times B')|\ge\theta|A'|b_0/4,
 \qquad |B'|\ge\theta b_0/4.
\]

Fix \(u\in A'\) and \(v\in B'\).  Of the at least
\(\theta|A'|/4\ge\theta|U|/8\) vertices in \(N(v)\cap A'\), at most
\(\theta|U|/16\) form a bad pair with \(u\).  Each remaining \(u_1\)
has at least \(\tau b_0\) common neighbors \(v_1\) with \(u\).
Thus \(u,v\) have at least \(c\theta^5a_0b_0\) original-graph walks
\(u,v_1,u_1,v\).  The walk determines
\[
 (u+v_1,\ u_1+v_1,\ u_1+v)\in C^3,
\]
and, for fixed \(u,v\), this map is injective.  Its alternating sum
is \(u+v\).  Choose one representing pair for every distinct element
of \(A'+B'\); the corresponding triples are disjoint.  Hence
\(c\theta^5a_0b_0|A'+B'|\le |C|^3\), as required.
\end{proof}

Here is the common finite-mesh input for the next three lemmas.  The
parameter law may be a law on indices: distinct direction bins with the
same rounded parameter are not silently identified.  The graph bounds
refer to the corresponding original edge sets.

\begin{assumption}[Planar graph data]\label{re:graph-data}
Let \(A,B\subset h\Z^2\cap[-K,K]^2\), with
\(h^{-1}/K\le |A|,|B|\le Kh^{-1}\).  A probability law \(\nu\) on
represented indices \(\kappa\) obeys, for a fixed \(S>0\),
\begin{equation}\label{re:ratio-cap}
 K^{-1}\le|\kappa|\le K,\qquad
 \nu(I)\le K|I|^S\quad(h\le |I|\le1).
\end{equation}
For each index there is \(G_\kappa\subset A\times B\) with
\begin{equation}\label{re:graph-hypotheses}
 |G_\kappa|\ge h^{-2}/K,\qquad
 N_h\{x+\kappa y:(x,y)\in G_\kappa\}\le Kh^{-1}.
\end{equation}
The uniform law on \(G_\kappa\) obeys, for intervals \(I,I'\) of
length \(h\),
\begin{equation}\label{re:first-coordinate-cap}
 \Pp(x_1\in I,y_1\in I')\le Kh^2.
\end{equation}
\end{assumption}

Set \(L=h^{-1}\) in this section.  For each \(\kappa\), round
\(A\) and \(\kappa B\) to \(h\Z^2\).  Since \(|\kappa|\) and its
reciprocal are subpower, a rounded vertex has at most \(K\)
preimages, after increasing \(K\); an image edge has at most \(K^2\)
preimages.  The image graph therefore has subpower density, while
the set of its edge sums has at most \(KL\) rounded values.
Apply \Cref{re:graph-bsg} to this image graph and pull back all
preimages of the selected vertices.  Every retained image edge has
an original preimage, so the result is not merely a vertex-set
statement.  We obtain \(A_\kappa\subset A\), \(B_\kappa\subset B\)
such that
\begin{equation}\label{re:individual-bsg}
 |A_\kappa|,|B_\kappa|\ge L/K,\qquad
 N_h(A_\kappa+\kappa B_\kappa)\le KL,\qquad
 |G_\kappa\cap(A_\kappa\times B_\kappa)|\ge L^2/K.
\end{equation}

We use the covering-number forms of the Ruzsa triangle and fixed-iterate
Pl\"unnecke--Ruzsa inequalities \citep{TaoVu2006}:
\begin{align}
 N_h(U-W)N_h(V)&\le C N_h(U-V)N_h(V-W),
 \label{re:ruzsa-triangle}\\
 N_h(mU-nU)&\le C_{m,n}K^{C_{m,n}}N_h(U)
 \quad\text{if }N_h(U-U)\le K N_h(U).
 \label{re:iterate}
\end{align}
These apply to \(h\)-cell unions after rounding, with only fixed
dimensional changes of mesh.  For the first inequality, choose one
representing pair for each separated output \(u-w\) and let \(v\)
range over a separated subset of \(V\); the rounded pair
\((u-v,v-w)\) determines both the output and \(v\), up to bounded
choices.  We also use the elementary covering lemma
\begin{equation}\label{re:covering}
 N_h(U+V)\le K N_h(V)
 \quad\Longrightarrow\quad
 U\subset\bigcup_{i=1}^{CK}(u_i+V-V+B(0,Ch)).
\end{equation}
To see this, take a maximal disjoint collection of translates of
\(V+B(0,h)\) inside \(U+V+B(0,h)\); volume comparison bounds their
number and maximality gives the inclusion.

The sets in \eqref{re:individual-bsg} depend on the parameter.  We next
choose one pair of endpoint sets that has large overlap with the sets
for many parameters.  The Ruzsa inequalities then control sums of
fixed numbers of dilates of this one pair, including the scalar
expressions needed later in the collision argument.

\begin{lemma}[A common pair and scalar closure]\label{re:common-sets}
After restricting \(\nu\) to a subpower fraction of its mass, there
are a represented \(\kappa_0\), fixed \(A_*\subset A\), \(B_*\subset B\),
and \(\mathcal R=\{\kappa/\kappa_0\}\) with the following properties.
At least \(L^2/K\) original edges of \(G_{\kappa_0}\) lie in
\(A_*\times B_*\).  Put \(E=A_*\), \(F=-\kappa_0B_*\).
For every fixed \(m,H\), the set
\begin{equation}\label{re:scalar-support}
 c_1E+\cdots+c_mE+d_1F+\cdots+d_mF
\end{equation}
has at most \(K_{m,H}L\) \(h\)-cells whenever each coefficient is a
sum of at most \(H\) signed products of at most \(H\) elements of
\(\mathcal R\).  A common subpower dilation of all coefficients is
permitted.  The restricted probability law on \(\mathcal R\) still
satisfies a bound of the form \eqref{re:ratio-cap}, with a new
subpower constant and the same fixed exponent \(S\).
\end{lemma}
\begin{proof}
The rectangles \(A_\kappa\times B_\kappa\) from
\eqref{re:individual-bsg} sit in the one universe \(A\times B\),
which has at most \(K^2L^2\) elements.  Independent indices satisfy
\[
 \E_{\kappa,\kappa'}
 |(A_\kappa\times B_\kappa)\cap
   (A_{\kappa'}\times B_{\kappa'})|
 \ge\frac{(\E_\kappa|A_\kappa\times B_\kappa|)^2}{|A\times B|}
 \ge L^2/K^6.
\]
Choose \(\kappa_0\) with at least this average intersection, then
retain indices with intersection at least \(L^2/(2K^6)\).  Their
probability is at least \(1/(2K^8)\).  Set
\(A_*=A_{\kappa_0}\), \(B_*=B_{\kappa_0}\).  Because each
individual set has at most \(KL\) points, the retained indices have
\[
 |A_\kappa\cap A_*|,\ |B_\kappa\cap B_*|\ge L/(2K^7).
\]
The edge assertion follows from \eqref{re:individual-bsg} at
\(\kappa_0\).  Restricting and renormalizing the index law costs only
a subpower factor in its interval bound; division by
\(\kappa_0\), whose norm and reciprocal are subpower, does likewise.

For clarity about the uniform set calculus, put
\[
 d_h(U,V)=\log\frac{N_h(U-V)}{\sqrt{N_h(U)N_h(V)}}.
\]
The triangle inequality \eqref{re:ruzsa-triangle} is the triangle
inequality for \(d_h\), up to \(O(1)\).  Changing mesh from \(h\)
to \(h/K^C\), or dilating both sets by a scalar with norm and
reciprocal at most \(K^C\), changes the bounds below by \(O_C(\log K)\):
a planar cell refines into at most \(O(K^{2C})\) cells.
The small cross sum in \eqref{re:individual-bsg} gives
\(d_h(A_\kappa,-\kappa B_\kappa)=O(\log K)\).  Applying the triangle
inequality in both directions gives small self-differences for
\(A_\kappa\) and \(B_\kappa\).  To use the large intersections,
put \(H_\kappa=A_\kappa\cap A_*\).  Both
\(A_\kappa-H_\kappa\) and \(H_\kappa-A_*\) are contained
in the corresponding self-difference sets.  Since
\(|H_\kappa|\ge L/(2K^7)\), the triangle inequality through
\(H_\kappa\) gives \(d_h(A_\kappa,A_*)=O(\log K)\).
The same argument for the intersections of the \(B\)-sets gives
\(d_h(B_\kappa,B_*)=O(\log K)\).  Consequently
\begin{equation}\label{re:fixed-cross}
 d_h(E,F)=O(\log K),\qquad
 d_h(E,tF)=O(\log K)\quad(t\in\mathcal R).
\end{equation}
Here the second relation uses \(tF=-\kappa B_*\).
To transfer this relation from \(F\) to \(E\), use the triangle
inequality through \(tF\) and then undo the common dilation:
\begin{align*}
 d_h(E,tE)
 &\le d_h(E,tF)+d_h(tF,tE)+O(1)\\
 &=d_h(E,tF)+d_h(F,E)+O(\log K)
 =O(\log K).
\end{align*}
The dilation error is uniform because both \(|t|\) and
\(|t|^{-1}\) are subpower bounded.  For
\(p_j=t_1\cdots t_j\), \(t_i\in\mathcal R\), put \(p_0=1\).
For each fixed \(j\), the same argument gives the successive bounds
\begin{align*}
 d_h(E,p_jE)
 &\le d_h(E,p_{j-1}E)
       +d_h(p_{j-1}E,p_{j-1}t_jE)+O(1)\\
 &\le d_h(E,p_{j-1}E)+d_h(E,t_jE)+O_j(\log K)
 =O_j(\log K),\\
 d_h(E,p_jF)
 &\le d_h(E,p_jE)+d_h(p_jE,p_jF)+O(1)\\
 &\le d_h(E,p_jE)+d_h(E,F)+O_j(\log K)
 =O_j(\log K).
\end{align*}
Here the first bound is inductive, starting with the small
self-difference of \(E\).  The norm and reciprocal of every
\(p_j\) remain subpower bounded when \(j\) is fixed.

To include signed products, \eqref{re:iterate} with \(m=2,n=0\)
first gives \(N_h(E+E)\le KL\), and hence
\(d_h(E,-E)=O(\log K)\).  For \(X=E\) or \(F\),
\begin{align*}
 d_h(E,-p_jX)
 &\le d_h(E,-E)+d_h(-E,-p_jX)+O(1)\\
 &=d_h(E,-E)+d_h(E,p_jX)+O(1)
 =O_j(\log K).
\end{align*}

More concretely, for every signed product dilate \(U\) of \(E\) or
\(F\), of bounded length, \(N_h(U-E)\le K_HL\), while
\(N_h(E)\ge L/K\).  Apply \eqref{re:covering} with \(-E\): \(U\)
lies in at most \(K_H\) translates of \(E-E+B(0,K_Hh)\).
A fixed sum of such dilates is therefore covered by \(K_{m,H}\)
translates of a fixed iterate of \(E-E\), with at most
\(K_{m,H}L\) cells by \eqref{re:iterate}.  For a coefficient
\(c=\sum_{\ell=1}^q\epsilon_\ell\alpha_\ell\), where
\(q\le H\), \(\epsilon_\ell\in\{-1,1\}\), and each
\(\alpha_\ell\) is a product of at most \(H\) elements of
\(\mathcal R\), we use only the containment
\[
 cX\subseteq \epsilon_1\alpha_1X+\cdots+\epsilon_q\alpha_qX,
 \qquad X=E\ \text{or}\ F.
\]
Indeed the left side uses the same point of \(X\) in every summand
on the right.  Thus each of the \(2m\) coefficient dilates in
\eqref{re:scalar-support} is contained in a sum of at most \(H\)
signed product dilates.  No sum of products is divided by, so
cancellation, including \(c=0\), does not affect the argument.
This proves \eqref{re:scalar-support}.
A common subpower dilation changes the cover by only a subpower
factor.
\end{proof}

\subsection{Flat scalar coefficients}

We now seek random coefficients with two properties: their values
belong to the scalar class just constructed, and their probabilities
on intervals of length \(h\) are nearly as small as \(h\).  The
first property keeps every fixed-coefficient endpoint sum in only
\(h^{-1-o(1)}\) mesh cells, forcing two independent samples to
collide with probability at least \(h^{1+o(1)}\).  The second will
give a smaller collision probability if two endpoint-difference
vectors are sufficiently transverse.  Their comparison in
\Cref{re:parallel-strips} will force the parallel-strip conclusion.

The analytic input is a finite-scale Fourier-decay theorem.  For a
probability measure \(\lambda\), write
\[
 I_s(\lambda)=\iint |x-y|^{-s}\,d\lambda(x)\,d\lambda(y),
 \qquad I_s^r(\lambda)=I_s(\lambda*P_r),
\]
where \(P_r\) is a fixed smooth approximate identity at scale
\(r\).  Theorem 1.5 of \citet{OrponenDeSaxceShmerkin} says: for
fixed \(n\ge2\) and \(s_j\in(0,1]\) with \(\sum_j s_j>1\), there are
\(r_0,\epsilon_0,\tau>0\), depending only on these fixed data, such
that Borel probability measures \(\lambda_j\) on \([-1,1]\) satisfying
\(I_{s_j}^r(\lambda_j)\le r^{-\epsilon_0}\), \(0<r<r_0\), obey
\begin{equation}\label{re:oss-input}
 |(\lambda_1\boxtimes\cdots\boxtimes\lambda_n)^\wedge(\xi)|
 \le |\xi|^{-\tau}\qquad(r^{-1}\le|\xi|\le2r^{-1}),
\end{equation}
where \(\boxtimes\) is multiplication of independent real variables.
The theorem does not require their supports to avoid zero.
The application below uses only exponents strictly below \(1\).

\begin{lemma}[Flat generated coefficients]\label{re:flat-coefficients}
Let \(\nu_{\mathcal R}\) be the restricted ratio law in
\Cref{re:common-sets}.  For every fixed \(e>0\), there are fixed
positive integers \(n,k\) and a subpower \(R\ge1\) such that, with
independent \(T_{ij}\sim\nu_{\mathcal R}\),
\begin{equation}\label{re:generated-coefficient}
 C=\sum_{i=1}^k\prod_{j=1}^n(T_{ij}/R)
\end{equation}
satisfies
\begin{equation}\label{re:coefficient-flatness}
 \sup_x\Pp(C\in[x-h,x+h])\le C_eh^{1-e}.
\end{equation}
Every value in its support is an expression allowed in
\eqref{re:scalar-support}, with the common subpower dilation
\(R^{-n}\).  Independent copies of \(C\) may therefore be used
simultaneously as coefficients in that cover bound.
\end{lemma}
\begin{proof}
Take \(R\ge\sup|T|\), with \(R\) and its reciprocal subpower, and
let \(\lambda\) be the law of \(T/R\) on \([-1,1]\).  Rescaling
\eqref{re:ratio-cap}, also for intervals of length above \(1/R\)
by the trivial bound, gives
\(\lambda(B(x,u))\le K' u^S\) for \(h\le u\le1\), with
\(K'=h^{-o(1)}\).  Fix \(0<s<\min(S,1)\) and \(n\ge2\) with \(ns>1\).
For \(h\le r\le1\), smoothing at scale \(r\) bounds the energy kernel
by \(C_s\max(r,|x-y|)^{-s}\).  The \(r\)-ball and dyadic annuli,
using the interval cap, give
\begin{equation}\label{re:energy-bound}
 I_s^r(\lambda)\le C_sK'
 \left(r^{S-s}+\sum_{j:2^jr\le1}(2^jr)^{S-s}+1\right)
 \le C_s'K'.
\end{equation}
Fix \(0<e_0<\min(e/2,1)\).  The constants \(r_0,\epsilon_0,\tau\)
from \eqref{re:oss-input} are now fixed.  For small enough \(h\),
all \(r\in[h,h^{e_0}]\) lie below \(r_0\), and subpower \(K'\)
makes \eqref{re:energy-bound} at most \(r^{-\epsilon_0}\)
uniformly over that interval.  Applying \eqref{re:oss-input} with
\(r=|\xi|^{-1}\) on successive annuli gives, for
\(\omega=\lambda^{\boxtimes n}\),
\[
 |\widehat\omega(\xi)|\le|\xi|^{-\tau}
 \qquad(h^{-e_0}\le|\xi|\le h^{-1}),
\]
up to an inessential fixed adjustment at the lower endpoint.
Choose fixed \(k\) with \(k\tau>2\), and put
\(\eta=\omega^{*k}\), the law of \eqref{re:generated-coefficient}.
Choose a nonnegative integrable \(\psi\ge1\) on \([-1,1]\) whose
Fourier transform is supported in \([-1,1]\), for example a suitably
dilated squared-sinc majorant.  Fourier inversion yields, uniformly
in \(x\),
\begin{align*}
 \eta([x-h,x+h])
 &\le \int\psi((u-x)/h)\,d\eta(u)\\
 &\le Ch\int_{|\xi|\le h^{-1}}|\widehat\omega(\xi)|^k\,d\xi\\
 &\le Ch\left(h^{-e_0}+
       \int_{h^{-e_0}}^{h^{-1}}t^{-k\tau}\,dt\right)
 \le Ch^{1-e_0}\le C_eh^{1-e}.
\end{align*}
The integers \(n,k\) are fixed for each \(e\), before \(h\to0\),
and the same \(R^{-n}\) multiplies every generated coefficient.
\end{proof}

\subsection{Parallel strips from collision probabilities}

\begin{lemma}[Weighted parallel strips]\label{re:parallel-strips}
Let \(E,F\) be the fixed sets in \Cref{re:common-sets} and let
\(\mu\) be the uniform law on a subset of at least \(h^{-2}/K\)
original edges of \(G_{\kappa_0}\cap(A_*\times B_*)\), after
the change \(y\mapsto-\kappa_0y\).  Suppose the first-coordinate
rectangle cap \eqref{re:first-coordinate-cap} holds for \(\mu\),
with a possibly enlarged subpower \(K\).  There are two parallel
strips, one in each endpoint plane, of width \(h^{1-o(1)}\),
containing \(h^{o(1)}\) of the \(\mu\)-mass.  Their common slope is
\(h^{-o(1)}\)-bounded.
\end{lemma}
\begin{proof}
We prove that a fixed-power failure of strip concentration would
force most sampled endpoint differences to include a transverse
pair.  The generated coefficients then give incompatible lower
and upper bounds for the probability that two endpoint sums
occupy the same mesh cell.
Let \(D\le K\) bound the endpoint norms and define
\[
 p(w)=\sup\{\mu(S_E\times S_F):
       S_E,S_F\text{ parallel strips of width }w\}.
\]
Fix \(0<\varepsilon<1\) and \(c>0\).  Suppose along a sequence
that \(p(w)\le h^c\) for \(w=h^{1-\varepsilon}\).
Draw \(m\) independent pairs of independent edges
\((a_i,b_i),(a_i',b_i')\sim\mu\), and put
\[
 v_i=a_i-a_i',\qquad u_i=b_i-b_i',\qquad
 q=h^{\varepsilon/2},\qquad \Delta=qw/10.
\]
Summing the rectangle cap over the \(O(D/h)\) possible cells for
the other first coordinate gives a marginal cap
\(\Pp((a_i)_1\in I)\le CKD h\) for an \(h\)-interval \(I\).
Independence of the two sampled edges therefore gives
\(\Pp(|v_i|<q)\le CKD(q+h)\).
We claim
\begin{equation}\label{re:bad-wedge}
 \Pp\left(\max_{r,s\in\{v_i,u_i:1\le i\le m\}}
                   |r\wedge s|\le\Delta\right)
 \le[CKD(q+h)]^m+m p(w)^{m-1}.
\end{equation}
The first term covers the event that every \(|v_i|<q\).
Otherwise choose an index \(i\) with \(|v_i|\ge q\).  Small wedges
put every other \(v_j,u_j\) within width \(w/2\) of the line
\(\R v_i\).  Fix the anchor edges and all primed edges.
For each \(j\ne i\), its unprimed edge must then lie in one pair
of parallel strips of width \(w\), an event of probability at most
\(p(w)\).  These unprimed edges are independent.  A union bound over
the anchor proves \eqref{re:bad-wedge}.  Choose \(m\) fixed so large
that \(m\varepsilon/2>3\) and \(c(m-1)>3\); subpower \(K,D\) then make
the right side at most \(h^3\) for sufficiently small \(h\).

Fix \(e<\varepsilon/8\).  Independently draw \(C_i,D_i\), \(1\le i\le m\),
from the law in \Cref{re:flat-coefficients}.
Use the \emph{same coefficient tuple} in two independent endpoint sums:
\[
 Z=\sum_{i=1}^m(C_i a_i+D_i b_i),\qquad
 Z'=\sum_{i=1}^m(C_i a_i'+D_i b_i').
\]
For every fixed coefficient tuple, \Cref{re:common-sets} bounds
the support of each sum by \(K_mh^{-1}\) cells.  Conditional on the
tuple, \(Z,Z'\) are independent samples from the same distribution.
Cauchy--Schwarz on those cells therefore gives
\begin{equation}\label{re:collision-lower}
 \Pp(|Z-Z'|_\infty\le h\mid(C_i,D_i)_{i=1}^m)
 \ge h/K_m.
\end{equation}

Averaging over the coefficient tuple gives the same unconditional
lower bound.  For the upper bound, first fix the sampled endpoints.
Outside the event in \eqref{re:bad-wedge}, choose two difference
vectors \(r,s\) from the list \(\{v_i,u_i\}\) with
\(|r\wedge s|>\Delta\), making the choice from endpoints alone.
Condition on every coefficient except the two attached to \(r,s\).
They are still independent copies of the flat coefficient law,
even if the two vectors have the same index \(i\).  The condition
\(|Z-Z'|_\infty\le h\) restricts these two real coefficients to a
parallelogram of area \(O(h^2/\Delta)\) and perimeter
\(O(Dh/\Delta)\).  It meets at most
\(C((1+D)/\Delta+1)\) squares of side \(h\): count interior squares
by area and boundary squares by perimeter.  Independence and
\eqref{re:coefficient-flatness} give probability at most
\(C_e^2h^{2-2e}\) for each square.  Hence
\begin{equation}\label{re:collision-upper}
 \Pp(|Z-Z'|_\infty\le h)
 \le h^3+K_m\frac{h^{2-2e}}{\Delta}
 =h^3+K_mh^{1+\varepsilon/2-2e}.
\end{equation}
This contradicts \eqref{re:collision-lower}, since
\(\varepsilon/2-2e>0\) and \(K_m\) is subpower.

Thus for every fixed \(\varepsilon,c>0\), eventually
\(p(h^{1-\varepsilon})>h^c\).  Take a sufficiently slow diagonal
\(\varepsilon,c\downarrow0\), after the fixed integers and their
thresholds at each stage have been chosen.  This gives width
\(w=h^{1-o(1)}\) and mass \(m_0=h^{o(1)}\); no expression of
unbounded length is used at any fixed tolerance.
If \(q_0\) is a unit vector along the common strip direction,
the first-coordinate projection of either strip within \([-D,D]^2\)
has length at most \(C(D|(q_0)_1|+w)\).  Summing
\eqref{re:first-coordinate-cap} over the resulting pairs of
\(h\)-intervals gives
\[
 m_0\le CK(D|(q_0)_1|+w+h)^2.
\]
Since \(m_0\) is subpower below and \(K,D\) subpower above, while
\(w\to0\) by a fixed power at each diagonal stage,
\(|(q_0)_1|\ge h^{o(1)}\).  The common slope
\((q_0)_2/(q_0)_1\) is therefore subpower bounded.
\end{proof}

\begin{proposition}[Retained-edge planar rigidity]\label{re:retained-edge-rigidity}
Under \Cref{re:graph-data}, after increasing to a new subpower
\(K_*\), there are a represented \(\kappa_0\), at least
\(h^{-2}/K_*\) \emph{original} edges of \(G_{\kappa_0}\), and
\(|\lambda|+|c_A|+|c_B|\le K_*\) such that on those edges
\[
 |x_2-\lambda x_1-c_A|+
 |y_2-\lambda y_1-c_B|\le K_*h.
\]
No fixed-power dependence of \(K_*\) on the initial subpower
\(K\) is asserted.
\end{proposition}
\begin{proof}
The rounding and \Cref{re:graph-bsg} give \eqref{re:individual-bsg}.
\Cref{re:common-sets} fixes \(\kappa_0,A_*,B_*\) while retaining
\(h^{-2}/K\) edges of \(G_{\kappa_0}\); its scalar closure and
\Cref{re:flat-coefficients} meet the inputs of
\Cref{re:parallel-strips}.  Give those retained original edges
their uniform law.  Conditioning on a subset of size at least
\(h^{-2}/K\) inflates \eqref{re:first-coordinate-cap} only by a
subpower factor, since \(|G_{\kappa_0}|\le K^2h^{-2}\).
The invertible dilation \(y\mapsto-\kappa_0y\) likewise changes
an \(h\)-interval into at most \(K\) \(h\)-intervals, preserving
the cap with another subpower factor.  \Cref{re:parallel-strips}
selects subpower mass of these same edges.  Undoing that dilation
preserves the common slope and edge identities; it changes strip
width and intercept by at most a subpower factor.  Endpoint norms
bound both intercepts once the slope is subpower bounded.
Absorb all losses, including the slow diagonal, into \(K_*\).
\end{proof}

\subsection{An affine--skew row from line incidences}

We apply the planar rigidity theorem to a row evaluated along moving
lines.  In a fixed direction, the line intercept and the row's
evaluation are nearly constant.  Long line segments therefore give
small planar projection sets.  Two directions provide the endpoints
of a graph, and a third direction provides its small restricted
sumset.

\begin{lemma}[Affine--skew row fit]\label{a06:row-fit}
Let \(z(t)=(t,Z(t))\), \(0\le t\le1\), be affine trajectories
with velocity \((1,v)\), and suppose positions and velocities are
bounded by a subpower factor \(K\).  Let \(\mu\) be an incidence
measure of mass at least \(1/K\), dominated by \(K\) times a
trajectory prior times Lebesgue measure in time.  At a fine fixed-power
mesh \(\xi\), suppose a probability law \(\pi\) on velocities
satisfies \(\pi(I)\le K|I|^S\) for some fixed \(S>0\) and all
intervals with \(\xi\le|I|\le1\).  Suppose also that every pair
of position and velocity cells obeys
\[
 \mu(z_\xi,v_\xi)\le K\xi^2\pi(v_\xi).
\]
Let \(T\in[-K,K]^2\) be a row on the incidences.  Suppose it has
at most \(K\) bins of width \(E\ge\xi\) per position cell
\(z_\xi\), and \(T(1,v)\) lies within \(KE\) of a constant
depending only on the trajectory.  Then a restriction of mass at
least \(1/K\) satisfies
\[
 T=B+\lambda Jz+O(KE),\qquad
 J(t,Z)=(-Z,t),\qquad |B|+|\lambda|\le K.
\]
The subpower factor \(K\) may increase in the conclusion.
\end{lemma}
\begin{proof}
We may assume \(\xi\le E\le1\), since the conclusion is
immediate for \(E>1\).  Put \(V_0=(\xi/E)T\), and consider
the occupied cells of \((z,V_0)_\xi\).  There are at most
\(K\) per position cell and at most \(K\xi^{-2}\) in all.
For a slope bin \(d\) of width \(\xi\), with representative
\(v_d\), put \({\bf u}_d=(1,v_d)\).  On point-cell centers
use the linear projection
\[
 \mathcal P_{{\bf u}_d}(z,V_0)
    =(z\cdot J{\bf u}_d,\ V_0\cdot{\bf u}_d).
\]
For the true direction, its first coordinate is constant on a line;
the second is constant to error \(K\xi\) by the row hypothesis.

\paragraph{Projection covers and point-mass witnesses.}
Freeze the present incidence law as \(\mu_{\rm cov}\).
Delete lines with too little labeled time or incidence weight.
Prior-times-time domination leaves each used trajectory with a
\(\mu_{\rm cov}\)-witness time set of length at least \(1/K\).
It visits at least \(1/(K\xi)\) witness point cells, throughout
which its projection varies by at most \(K\xi\).  Replacing its
true velocity by the bin representative costs only \(K\xi\)
in both projection coordinates.

For each direction bin, projected values separated by a sufficiently
large multiple of \(K\xi\) consequently use disjoint witness
sets.  Since the total number of point cells is at most
\(K\xi^{-2}\), its currently occurring projection can be
covered by at most \(K\xi^{-1}\) mesh cells.  The law
\(\mu_{\rm cov}\) will continue to supply these covers, even
after the incidence law is restricted further.

Discard slope bins of palette weight below \(\xi^2\).
The joint cap, summed over at most \(K/\xi\) such bins and
the position cells, makes their total mass negligible.
Pigeonhole comparable palette weights on dense incidence mass,
leaving \(M\) bins.  Their total palette weight is at least
\(1/K\).  The uniform law on them therefore has an interval
bound \(K\rho^{S'}\), for some fixed \(S'>0\) and
\(\xi\le\rho\le1\).

Normalize this current incidence law to probability and freeze it
as \(\mu_{\rm pt}\).  Every joint point-direction cell has
mass at most \(K\xi^2/M\).  Delete point cells of marginal
mass below \(\xi^2/K_1\), and point-direction edges of mass
below \(\xi^2/(K_1M)\), with a sufficiently large subpower
\(K_1\).  The numbers of cells and edges pay for these
deletions.  At least \(\xi^{-2}M/K\) unweighted edges remain.
Deleting low-degree points leaves at least \(\xi^{-2}/K\)
point cells of degree at least \(M/K\).

Every surviving point still has its floor in the frozen
\(\mu_{\rm pt}\)-marginal.  Thus a later selection of
\(\xi^{-2}/K\) such point cells will lift to dense mass by
restoring all their \(\mu_{\rm pt}\)-incidences, regardless
of which graph directions were selected.  This second frozen law
supplies point-mass floors; the first supplies the long-line
projection covers.

\paragraph{Two projections and the third-direction sumset.}
At each surviving point there are at least \(M^3/K\) ordered
triples of incident slopes separated pairwise by at least
\(1/K\).  Indeed the interval bound makes the fraction within
a sufficiently small reciprocal-subpower neighborhood negligible
compared with the degree fraction.  Pigeonhole the first two
directions \(v_1,v_2\).  They have at least
\(\xi^{-2}M/K\) common-neighbor incidences with third
directions also separated from them.  Retain third directions with
at least \(\xi^{-2}/K\) such neighbors.  There are at least
\(M/K\) retained directions; use their uniform law.

Let \(\mathcal A\) and \(\mathcal B\) be the projections of
the common-neighbor point cells under
\(\mathcal P_{{\bf u}_1}\) and \(\mathcal P_{{\bf u}_2}\),
rounded to the \(\xi\)-lattice.  The earlier projection covers
give \(|\mathcal A|,|\mathcal B|\le K\xi^{-1}\).
Each endpoint pair has at most \(K\) point preimages, since
the basis determinant is at least \(1/K\) and its entries are
bounded by \(K\).  For each retained third direction, its
common-neighbor points therefore define a graph
\(G_\kappa\subset\mathcal A\times\mathcal B\) with at least
\(\xi^{-2}/K\) distinct edges.

For that direction write
\[
 {\bf u}_d=c_1{\bf u}_1+c_2{\bf u}_2,\qquad
 c_1=\frac{v_2-v_d}{v_2-v_1},\qquad
 c_2=\frac{v_d-v_1}{v_2-v_1},\qquad
 \kappa=\frac{c_2}{c_1}=\frac{v_d-v_1}{v_2-v_d}.
\]
The map \(\mathcal P_{\bf u}\) is linear in \({\bf u}\).
Thus, for the unrounded projections \(a,b\) of a common point,
\[
 a+\kappa b=c_1^{-1}\mathcal P_{{\bf u}_d}(z,V_0).
\]
Both \(c_i\) and their reciprocals are subpower bounded.
Rounding and this dilation cost only \(K\), so the third-direction
projection cover gives
\[
 N_\xi\{a+\kappa b:(a,b)\in G_\kappa\}
       \le K\xi^{-1}.
\]

We verify the remaining planar hypotheses.  The graph sizes imply
\(|\mathcal A|,|\mathcal B|\ge\xi^{-1}/K\), and both endpoint
sets are bounded by \(K\).  Separation from the basis slopes
gives \(K^{-1}\le|\kappa|\le K\).  The ratio law has the
same positive-exponent interval bound up to \(K\), since
\[
 |v_d-v_{d'}|\le K|\kappa_d-\kappa_{d'}|
\]
on the retained choices.  Finally, prescribing length-\(\xi\)
intervals for the first endpoint coordinates \((a_1,b_1)\)
restricts \(z\) to at most \(K\) cells by the invertible
basis map.  There are only \(K\) point options above those
cells.  Each graph therefore has only \(K\) possible edges
at such a first-coordinate rectangle.  Its uniform edge law
obeys the required cap, as does the uniform law on any subset
of at least \(\xi^{-2}/K\) edges.  Every such subset uses
at least \(\xi^{-2}/K\) distinct surviving point cells.

\paragraph{Recovering the row and restoring incidence mass.}
Apply \Cref{re:retained-edge-rigidity} with \(h=\xi\).
For a represented \(\kappa_0\), it retains at least
\(\xi^{-2}/K_*\) original edges of \(G_{\kappa_0}\) in
parallel strips of width \(K_*\xi\), with common slope and
intercepts bounded by a new subpower \(K_*\).  No fixed-power
dependence on the earlier \(K\) is needed.  Absorb \(K_*\)
into the subsequent subpower factor.

On the corresponding point cells,
\[
 V_0\cdot{\bf u}_i
   =b_i+\lambda_0z\cdot J{\bf u}_i+O(K\xi),\qquad i=1,2.
\]
The estimate depends only on \((z,V_0)\) and remains valid
throughout each selected cell.  Inverting the basis and undoing
\(V_0=(\xi/E)T\) gives
\(T=B+\lambda Jz+O(KE)\) there.

Bounded endpoint multiplicity leaves at least \(\xi^{-2}/K\)
distinct selected point cells.  Each has frozen
\(\mu_{\rm pt}\)-mass at least \(\xi^2/K_1\), so their
whole-cell mass is dense.  Restore those incidences.  The row
estimate applies even when their directions were not among the
selected graph edges, and no incidence removed before freezing
\(\mu_{\rm pt}\) is restored.

It remains to bound the coefficients after undoing the scaling.
If \(|\lambda|\) were power-large, the bound \(|T|\le K\)
would confine this dense selected \(z\)-mass to a ball of
radius \(K/|\lambda|\).  Summing the joint cell cap over
velocities bounds its mass by
\[
 K(K/|\lambda|+\xi)^2,
\]
a contradiction.  Hence \(|\lambda|\le K\), and an occupied
point then bounds \(|B|\) by \(K\) as well.
\end{proof}
